\section{Experimental Setup}
\label{sec:setup}

\subsection{Datasets and leakage-free protocols}
\label{sec:data}
\textbf{RadHAR (MMActivity)} \cite{singh2019radhar}: TI IWR1443 radar, 30\,Hz,
five activities (boxing, jumping jacks, jumping, squats, walking), 200
recordings (71.5\,min train, 21.4\,min test), median 21 returns per frame. We
parse the raw ROS dumps and recover the official 71.6/21.4-minute partition. As
in \cite{singh2019radhar}, samples are 2-s windows (60 frames) with a 10-frame
stride (3{,}565 test windows; the original paper reports 3{,}538). We keep the
\emph{official test recordings} unchanged. The validation set, however, holds
out 20\% of the \emph{training recordings} per class (29 recordings) instead of
20\% of the windows. Neighbouring windows share 50 of 60 frames, so a
window-level split would leak near-duplicates into the data used for model
selection. A few training recordings are 120--190\,s long instead of 20\,s, so
each validation recording contributes at most 55 evenly spaced windows. This
keeps one long recording from dominating model selection (1{,}595 validation
windows). Training windows use a 20-frame stride (4{,}773 windows), the same
for every model.

\textbf{MiliPoint} \cite{cui2023milipoint}: TI IWR1843, 24\,Hz, 11 subjects, 49
actions. 212{,}920 frames carry action labels, matching the original paper. The
median frame has 6 returns, with $xyz$ only. Labelled frames are grouped into
1{,}139 contiguous single-action runs, and 2-s windows (48 frames) never cross a
run boundary. Our \textbf{primary protocol is subject-disjoint}: subjects 8--10
form the test set (3{,}569 windows), subjects 6--7 the validation set, and the six
remaining subjects the training set. For comparison with published numbers we
also evaluate (a) the \emph{chronological} 80/10/10 cut used by the MiliPoint
code and (b) a \emph{random} 80/10/10 split of overlapping windows, which is
leaky by construction.

\subsection{Baselines}
All models were re-implemented in pure PyTorch and trained under identical data,
augmentation and optimisation settings:
the \textbf{voxel TD-CNN+BiLSTM} of RadHAR \cite{singh2019radhar} (per-segment
$10\times32\times32$ occupancy grids, 3-D convolutions with 8/16/16 filters, a
reduced width for CPU training);
\textbf{PointNet} \cite{qi2017pointnet};
\textbf{PointNet+BiLSTM}, a frame-level temporal baseline in the spirit of
\cite{gu2024lpn};
\textbf{PointNet++} (single-scale grouping) \cite{qi2017pointnet2};
\textbf{DGCNN} \cite{wang2019dgcnn};
and \textbf{Point Transformer} \cite{zhao2021pt}.
The static set models receive all $G\!\times\!P$ points with time as a feature.
DGCNN and Point Transformer use a 128-point subset because their cost grows
quadratically with the number of points. We also compare against published
results under their own protocols (Table~\ref{tab:sota}).

\subsection{Training and evaluation details}
All models use AdamW \cite{loshchilov2019adamw} (learning rate
$2\!\times\!10^{-3}$, weight decay 0.05), a one-cycle schedule, batch size 64,
label smoothing 0.1 and gradient clipping at 1.0. Training runs for at most 15
epochs with early stopping (patience 5) on \emph{validation macro-F1}. The
selected checkpoint is evaluated once on the test set. Every configuration is
trained with three seeds (0, 1, 2), and we report mean $\pm$ sample standard
deviation. For the reference model we give 95\% bootstrap confidence intervals
(2{,}000 resamples of test windows) \cite{efron1979bootstrap}. Paired
comparisons use exact McNemar tests \cite{mcnemar1947} with Holm correction
\cite{holm1979} within each seed. Metrics are top-1 accuracy and macro-F1, plus
top-3 accuracy for MiliPoint. For efficiency we report parameters, FLOPs per
window, FP32/INT8 storage, and median CPU latency at batch size 1 on a single
thread (Intel Xeon, PyTorch~2.15). Single-thread CPU latency stands in for an
edge processor. All experiments ran on a 2-core CPU without a GPU.
